% SPDX-License-Identifier: MIT
\chapter{Layout, Fokus und Ereignisse}\label{chap:layout}
\ChapterSummary{Sie lernen die Zweiphasen-Messung kennen,
wählen geeignete Layout-Container und verstehen, wie Tastatur-,
Text- und Pointer-Ereignisse zum richtigen Widget gelangen.
Das Kapitel vermittelt die Grundlage für reaktionsfähige,
gut testbare Oberflächen.}

\section{Warum Layout in zwei Phasen erfolgt}
Ein Widget weiß vor der Darstellung nicht zwingend, wie viel
Platz seine Beschriftung benötigt: Ein Terminal misst in
Zeichen-Zellen, eine gerenderte Oberfläche in Schriftmetriken.
Darum verwendet das Toolkit \term{measure} und \term{arrange}.

\term{measure} fragt mit Minimal- und Maximalbedingungen
nach der gewünschten Größe. Die \term{MeasurementContext}
liefert passende backendbezogene Messregeln, ohne dass
ein Core-Widget das Backend kennen muss.
\term{arrange} legt danach die endgültige Position
und Größe innerhalb des zugewiesenen Rechtecks fest.

\begin{lstlisting}[language=SASDCpp]
sasd::ui::terminal::TerminalMeasurementContext metrics;
const Size viewport{80, 24};
(void)form.measure(metrics, {{0, 0}, viewport});
form.arrange({0, 0, viewport.width, viewport.height});
\end{lstlisting}
Der Ausschnitt setzt voraus, dass die Core-Typen aus
\texttt{sasd::ui} bekannt sind und \texttt{form}
bereits existiert und Sie den Terminal-Header
\filepath{sasd/ui/terminal/terminal_measurement_context.hpp}
eingebunden haben. Eine echte Host-Anwendung passt
die Werte nach einem Resize erneut an.

\section{Welcher Layouttyp passt?}

\noindent\begin{tabularx}{\linewidth}{@{}p{30mm}X@{}}
\toprule
Container & Praktische Verwendung\\\midrule
\term{VBox} & Beschriftete Schritte, einfache Formulare,
Aktionsstapel; \texttt{setSpacing}.\\
\term{HBox} & Horizontale Befehlsleisten und kleine
zusammengehörige Bedienelemente.\\
\term{GridLayout} & Regelmäßige Spalten/Raster mit
konfigurierbaren Zeilen- und Spaltenabstaenden.\\
\term{FormLayout} & Label-Feld-Paare; mit
\texttt{emplaceRow<WidgetTyp>(...)}.\\
\term{StackLayout} & Mehrere Widgets belegen denselben
lokalen Inhaltsbereich, etwa für Overlays.\\
\bottomrule
\end{tabularx}
Das sind zurzeit konkrete Container, kein
CSS-/HTML-Layoutsystem. \term{VBox} und \term{HBox}
führen keine magische automatische Flexbox-Logik ein,
die beliebige UI-Fälle bereits abdeckt.

\begin{lstlisting}[language=SASDCpp,caption={FormLayout für beschriftete Felder}]
auto& fields = window.emplace<FormLayout>();
fields.setColumnSpacing(2);
auto nameRow = fields.emplaceRow<TextField>("Benutzer:");
nameRow.field.setText("admin");
\end{lstlisting}
Für Überlagerungen benutzen Sie \term{StackLayout}
nur dann, wenn das erwartete Sichtbarkeits- und
Fokusverhalten ausdrücklich geregelt ist.
Ein Stack bedeutet noch keinen kompletten
modale-Dialoge-Manager.

\section{Fokus ist logischer Zustand}
\term{FocusManager} verwaltet genau den gültigen
logischen Empfänger der Tastatureingabe.
\term{FocusTraversal} bestimmt eine reproduzierbare
Reihenfolge durch den sichtbaren Widget-Baum.
\textbf{Tab} und \textbf{Shift+Tab} können so Fokus
wechseln, ohne dass jeder Button die gesamte
Oberfläche kennen muss.

Die Wirksamkeit des Fokus hängt von
\texttt{isFocusable()}, \texttt{isVisible()} und
\texttt{isEnabled()} ab. Wenn ein zuvor fokussiertes
Element unbrauchbar wird, darf kein veralteter
Fokuszeiger unkontrolliert weiterverwendet werden.
Für große Formulare empfiehlt es sich, zuerst
die Tab-Reihenfolge festzulegen und danach
Shortcuts und Popups einzubauen.

\section{KeyEvent ist kein eingegebener Text}
Ein Druck auf die Taste A ist ein
\term{KeyEvent}; die vom Betriebssystem bzw.
Terminal bestätigte Texteingabe ist ein
\term{TextInputEvent}. Beide sind nicht
gleichzusetzen. Tastaturbelegung, AltGr,
Modifikatoren oder Eingabemethoden können
ein andersartiges Zeichen erzeugen als ein
physischer Tastencode vermuten lässt.

Die \term{EventQueue} sammelt Ereignisse,
\term{Application::processEvents()} verarbeitet sie.
\term{EventDispatcher} kann Ereignisse mit
\emph{handled/ignored}-Semantik vom Ziel
zu übergeordneten Elementen weiterreichen.
Der konkrete Host entscheidet, wie das Ziel
ermittelt, wann Fokus umgeschaltet und
wie die nächste Darstellung angestossen wird.

\section{Maus und Pointer: Capture hat einen Besitzer}
\term{PointerRouter} führt geometrisches Hit-Testing,
Hover und explizites Pointer-Capture zusammen.
Beim Ziehen darf nicht plötzlich ein anderes
Widget die laufende Geste übernehmen, nur weil
die Maus das ursprüngliche Rechteck verlässt.
Ein Button besitzt semantische Armed-/Pressed-Zustände;
der Router verwaltet, welches Widget die physische
Geste beansprucht. Verliert ein Widget seine
Capture-Berechtigung, muss es seine Geste sauber
zurücksetzen.

Das Terminal kennt Pointer-Ereignisse nur dann,
wenn das verwendete Terminal das entsprechende
Protokoll liefert und der Host dessen Aktivierung
gewollt hat. Eine alternative reine Tastaturbedienung
muss deshalb erhalten bleiben.

\section{Invalidierung und Darstellung}
Eine Textänderung kann anderes \term{Layout} erfordern,
eine reine Farb- oder Pressed-Änderung meist nur
eine neue \term{Darstellung}.
Das Toolkit unterscheidet Mess-Invalidierung und
visuelle Invalidierung. Der \term{PresentationCoordinator}
organisiert danach die anstehenden Präsentationsupdates,
und ein \term{PresentationSink} erzeugt den Ziel-Frame.
Bei Oberflächenveränderungen dürfen alte Regionen
nicht versehentlich sichtbar bleiben; dafür existiert
konservatives Refresh-Verhalten.

\Important{Callbacks sind keine Einladung zu blockierendem
I/O im Eventloop. Netzwerkabfragen oder lang laufende
Administrationsoperationen sollten die Interaktion
nicht unnötig anhalten. Thread- und Worker-Integration
ist allerdings kein bereits fertiger universeller
Async-Runtime-Vertrag des Toolkits.}

\PracticalTip{Testen Sie eine Form zuerst bei grossem,
dann bei sehr kleinem Terminalfenster. Beobachten Sie
Fokus, abgeschnittenen Inhalt, die Reaktion auf Resize
und ob deaktivierte Widgets korrekt übersprungen werden.}
